\documentclass[12pt]{article}
\usepackage[utf8]{inputenc}
\usepackage[T1]{fontenc}
\usepackage{hyperref}
\title{Kaapor Language Resources for K-12 Curriculum Development}
\author{Compiled by Claude Sonnet 5 (claude-sonnet-5) for Andr\'as Kornai}
\date{}
\begin{document}
\maketitle

\section{Introduction}
Kaapor (Ka'apor, Urubu-Kaapor; ISO 639-3: urb) is a Tupi-Guarani language of Maranhão, Brazil. Wikipedia gives about 800 native speakers (2006). Tembé people also speak it as a second language. A village sign language also exists in the community.

\section{Primary Educational Resources}
\begin{itemize}
\item A doctoral thesis by Raimunda Benedita Cristina Caldas (UnB, 2009) presents a model Ka'apor-Portuguese bilingual dictionary. It responds to a community request for documentation with a view to teaching the language in village schools: \url{https://repositorio.unb.br/bitstream/10482/8867/1/2009_RaimundaBeneditaCristinaCaldas.pdf}
\end{itemize}

\section{Linguistic and Reference Materials}
\begin{itemize}
\item Wikipedia on the Ka'apor language: \url{https://en.wikipedia.org/wiki/Ka%CA%BCapor_language}
\item Endangered Languages Project entry: \url{https://endangeredlanguages.com/lang/3151}
\item Two Portuguese-language journal articles surfaced in the search (not examined): \url{https://periodicos.ufac.br/index.php/RFIR/article/download/3095/1887/} and \url{https://recima21.com.br/recima21/article/download/1068/861/7425}
\end{itemize}

\section{Religious and Community Resources}
\begin{itemize}
\item Povos Indígenas no Brasil (ISA) page on the Ka'apor people: \url{https://pib.socioambiental.org/en/Povo:Ka'apor}
\end{itemize}

\section{Digital and Audio Resources}
\begin{itemize}
\item Wikipedia on Ka'apor Sign Language, relevant for deaf education: \url{https://en.wikipedia.org/wiki/Ka%CA%BCapor_Sign_Language}
\end{itemize}

\section{Curriculum Development Context}
The main finding is the bilingual dictionary model, which was designed with village-school use in mind. Its publication status and any school use were not established. Current speaker numbers and any existing school programmes also remain unverified.
\end{document}